\section{术语表与工具清单}
\subsection{常见术语}
为了方便团队协作，Graham 推荐维护术语表：如 Plan-Act-Verify、Agent Chain、Observability Stream 等，统一语言可以减少误解。

\begin{itemize}
    \item Plan-Act-Verify：明确划分 Agent 的计划、执行与验证步骤。
    \item Agent Chain：把多个 Agent 组合成复杂任务的链式执行流程。
    \item Observability Stream：Agent 输出的日志/命令/测试结果的统一数据流。
\end{itemize}

\subsection{工具清单}
OpenHands 内置了多套工具：Terminal Runner、Browser Controller、Test Executor、Diff Reporter。部署时可以为每个保持单独的环境与权限，并记录使用情况。

\begin{table}[H]
\centering
\begin{tabular}{p{0.3\textwidth} p{0.6\textwidth}}
\toprule
工具 & 用途 \\
\midrule
Terminal Runner & 运行 shell 命令并记录 stdout/stderr。 \\
\midrule
Browser Controller & 进行页面交互、抓取数据并回写环境。 \\
\midrule
Test Executor & 运行 pytest、npm test 等并收集结果。 \\
\midrule
Diff Reporter & 汇总 patch、生成 markdown 描述与差异。 \\
\bottomrule
\end{tabular}
\caption{OpenHands 主要工具及其角色。}
\end{table}

\begin{knowledgebox}{工具组合策略}
尽量在不同的步骤中使用不同权限的工具，降低误操作风险；比如把数据库操作限制在单独的命令中，并设置回滚逻辑。
\end{knowledgebox}

\subsection{本章小结}
\begin{itemize}
    \item 术语表让跨团队合作更顺畅，避免对 Agent 意图的误解。
    \item 工具清单帮助把不同类型的操作模块化，便于设置权限与审计。
\end{itemize}

